% Slide 5: Related Works & Contributions
\begin{frame}{\ModuleFrameTitle{2}{Related Works \& Contributions}}
\begin{columns}[T,totalwidth=\textwidth]

% ---- Left column: related works --------------------------------------
\column{0.50\textwidth}

{\small\textbf{Related Works}}

\vspace{0.10cm}
{\small\begin{itemize}\setlength\itemsep{2pt}
  \item \textbf{Perpetual grid exploration} without obstacles by luminous robots
        {\footnotesize\textcolor{softgray}{[Bramas et al., Rauch et al.]}}
  \item \textbf{Exploration with obstacles} restricted to oblivious/non-myopic robots
        {\footnotesize\textcolor{softgray}{[Prencipe, Sardar et al.]}}
  \item \textbf{Formal verification} via proof assistants (Rocq/Pactole) \& model-checking
        {\footnotesize\textcolor{softgray}{[Bérard et al., Courtieu et al.]}}
  \item \textbf{Algorithm synthesis} limited to rings and few robots due to combinatorial explosion
        {\footnotesize\textcolor{softgray}{[Bonnet et al., Millet et al.]}}
\end{itemize}}

% ---- Right column: contributions -------------------------------------
\column{0.46\textwidth}

{\small\textbf{Our Contributions}}

\vspace{0.10cm}
{\small\begin{itemize}\setlength\itemsep{2pt}
  \item \textcolor{accent}{\textbf{First generic methodology}} for synthesizing perpetual exploration algorithms in grids containing an obstacle.
  \item \textcolor{accent}{\textbf{Partly automated pipeline}}: scenario-guided synthesis $\rightarrow$ small-grid simulations $\rightarrow$ induction proof.
  \item \textcolor{accent}{\textbf{New valid algorithms}}: 5 (visibility~1) and 115 (visibility~2) new algorithms extending existing baselines.
  \item \textcolor{accent}{\textbf{Automated classification}}: categorized by syntactic and performance metrics.
\end{itemize}}

\end{columns}
\end{frame}